\subsection{位数4のrank2基底の双対化と正常二次}
\label{subsec:cubic-fourth-root-unit-simple-rank2}
\begin{lemma}[rank2基底には非退化な正常二次がある]
\label{lem:cubic-fourth-root-rank2-nondegenerate-normal}
$\mu_{B,s}=(\rho,\rho,1,-1)$、$\rho=\pm i$ の必要atlasで偶射影rank2を仮定する。
$bc=h(U),cb=k(U)$、$S=\langle U\rangle$ とすると、$h,k$ の少なくとも一方は
$U^{\otimes2}$ のrank2 tensorである。
\end{lemma}
\begin{proof}
$h,k\ne0$、$S_2=U^{\otimes2}$、$h=f(k),k=g(h)$、
$h u=fg(u)h,k u=gf(u)k$ は既証のrank2構造による。
まず $h,k$ が比例するrank1 tensorなら、Frobenius middleの $U^2$ 対 $bc,cb$ の二列は
ともにrank1 tensorの倍数となる。
巡回で得るmixedの $2\times2$ 小blockもrank高々1となるが、
二つの可逆mixed graphと非退化pairingは四小block全てを可逆にする。
具体的にmixed全blockは $\binom{F}{I}M\binom{G}{I}^{t}$ の形で、
$F,G,M$ は可逆 $2\times2$ 行列である。
四小blockは $FMG^t,FM,MG^t,M$ となる。矛盾する。

従って両方rank1とすれば独立である。
$h=u v,k=s t$ と分解し、$f(k)=h,g(h)=k$ を自由な二次tensorで比較すると
$fg$ は $ku,kv$ をそれぞれ保存する。
$u,v$ が独立なら $h u=fg(u)h$ の左因子 $u$ を消去して
$vu=\lambda uv$ を得るが、$S_2$ が自由なことに反する。
よって $h$ は平方、$k$ も同様に平方であり、基底を選んで $h=x^2,k=y^2$ とできる。
$\mu_B(h)=-h,\mu_B(k)=-k$ と $S_2$ の自由性から
$\mu_B(x)=\rho x,\mu_B(y)=\rho y$。

even middleはrank2で、その関係基底 $bc-h,cb-k$ の $bc,cb$ 射影は同型である。
従ってその $U^2$ 主block $P$ はrank2で像が $\langle x^2,y^2\rangle\subset\operatorname{Sym}^2U$ となる。
対応する純 $U$ 四次tensor $w_U$ は第一・第二位置の交換に不変。
一方、中山巡回から $\tau(w_U)=-\rho w_U$。
位置交換を巡回で共役すると全隣接交換への不変性が従い、$w_U$ は完全対称となる。
すると $\tau(w_U)=w_U$ であり、$-\rho\ne1$ より $w_U=0$。
これは $P$ のrank2に反する。
\end{proof}

\begin{lemma}[同じAS indexを持つrank2基底の中山比較]
\label{lem:cubic-rank2-base-Nakayama-comparison}
上の基底 $S$ はnoetherian cubic AS正則次元3、AS index4であり、
\[
 \mu_S=\mu_B|_S.
\]
\end{lemma}
\begin{proof}
前補題の非退化な方を $h$ と書く。
$h$ は $S$ 上正規正則で、$H_S=1/((1-t)^2(1-t^2))$ から
$H_{S/(h)}=(1-t)^{-2}$。
非退化二次関係 $h$ のAS2平面が $S/(h)$ へ全射し、Hilbert級数の一致から同型である。
$S$ は既に左右noetherian、有限大域次元高々4である。
正常Rees lemmaと $S/(h)$ 上の長さ2の $k$ の解消を持ち上げることにより、
$S$ はAS Gorenstein次元3・index $2+2=4$、$\operatorname{pd}_S k\le3$ を得る。
従って $S$ はAS正則次元3である。

$L=B[b^{-1}]=S[b,b^{-1};f]$、$C=S[t;f]$、$\deg t=1$ とする。
$C$ の自己同型型Ore拡大の二項相対bimodule解消は
AS正則次元4・index5と $\mu_C|_S=f^{-1}\mu_S$ を与える。
$t$ を $b$ と同定すると $C[t^{-1}]=L$。
有限projective bimodule解消を左右flatなOre局所化で底変更すれば、
最高Extの二つの表示は
\[
 \underline{\operatorname{Ext}}^4_{L^e}(L,L^e)
 \simeq {}^1L^{\mu_B}(4)
 \simeq {}^1L^{\mu_C}(5)
\]
となる。
$L$ の単元は $\lambda b^n$ だけである。
実際、Laurent chargeの最大値と最小値は積で加法的なので単元は一つのchargeに限られ、
その係数は連結次数付き整域 $S$ の単元、従ってscalarである。
二つのshiftの生成元の次数は $-4,-5$ なので、上の同型は後者の生成元を
$\lambda b^{-1}$ 倍した前者へ送る。
右作用の両立性は
\[
 \mu_C(s)b^{-1}=b^{-1}\mu_B(s)\quad(s\in S)
\]
を与え、従って $\mu_C|_S=f^{-1}\mu_B|_S$。
Ore公式と比較して $\mu_S=\mu_B|_S$ を得る。
shift符号を省略して無次数のinner同値だけを比較した議論ではない。
\end{proof}

\begin{theorem}[単純固有値1を持つ位数4枝の全排除]
\label{thm:cubic-fourth-root-unit-simple-rank2-empty}
$(4;1,0,2)$ と $(4;3,0,2)$ のrank2必要atlasにはAS4整域が存在しない。
rank1排除と合わせ、この二スペクトルのH3候補は全て空である。
\end{theorem}
\begin{proof}
前二補題の非退化正常二次を $h$、そのnormalizerを $\tau$ とする。
$T=S/(h)$ は二元AS2平面である。
$\mu_S(h)=\mu_B(h)=-h$、$\tau(h)=h$ であり、$S_2=U^{\otimes2}$ なので、
これらは自由二次tensor上の厳密な等式である。
従って $T$ 上の自己同型として
\[
 \operatorname{hdet}_T(\mu_S)=-1,\qquad
 \operatorname{hdet}_T(\tau)=1.
\]
ここで二元AS2のhdetは定義二次関係のscalarである。
正常商の中山公式\cite[Lemma 1.5]{RRZ}は $\mu_T=(\mu_S\tau)|_T$ を与えるため、
$\operatorname{hdet}_T\mu_T=-1$ となる。
しかしAS2の量子・Jordan標準形で中山は関係をちょうど固定し、hdetは1である。
標数0で矛盾する。
\end{proof}

\begin{corollary}[三固有値領域の11個の必要スペクトル]
\label{cor:cubic-three-eigenvalue-11-list}
三固有値H3の必要リストは、位数3の3個、位数4の8個、合計11個に限られる。
各関係重みパターンは一つ、Galoisラベルは6軌道である。
残る候補の存在・H3許容性、四単純及びCY領域の全被覆は未分類である。
\end{corollary}
